\documentclass{article}
\usepackage{amsmath,amssymb}
\usepackage{xurl}
\usepackage[colorlinks=true,linkcolor=blue,citecolor=blue,urlcolor=blue]{hyperref}
% Shared commands for notes in docs/paper-gaps/.

\DeclareUnicodeCharacter{2080}{\ifmmode{}_{0}\else\textsubscript{0}\fi}
\DeclareUnicodeCharacter{2081}{\ifmmode{}_{1}\else\textsubscript{1}\fi}
\DeclareUnicodeCharacter{2082}{\ifmmode{}_{2}\else\textsubscript{2}\fi}
\DeclareUnicodeCharacter{2099}{\ifmmode{}_{n}\else\textsubscript{n}\fi}

\newcommand{\Lean}{\textsc{Lean}}
\ExplSyntaxOn
\NewDocumentCommand{\leanid}{m}{
  \ifmmode
    \textsf{\detokenize{#1}}
  \else
    \group_begin:
    \urlstyle{sf}
    \tl_set:Nn \l_tmpa_tl {#1}
    \tl_replace_all:Nnn \l_tmpa_tl {\_} {_}
    \exp_args:NV \path \l_tmpa_tl
    \group_end:
  \fi
}
\ExplSyntaxOff
\providecommand{\tr}{\operatorname{tr}}
\newcommand{\tnleanIssueBase}{https://github.com/LionSR/TNLean/issues/}
\newcommand{\tnleanPrBase}{https://github.com/LionSR/TNLean/pull/}
\newcommand{\tnleanDocBase}{https://sirui-lu.com/TNLean/docs/}
\newcommand{\ghissue}[1]{\href{\tnleanIssueBase#1}{GitHub issue \##1}}
\newcommand{\ghpr}[1]{\href{\tnleanPrBase#1}{PR \##1}}
\newcommand{\doclink}[1]{\href{\tnleanDocBase#1.html}{\path{#1}}}

% Dirac notation and the identity operator, as in blueprint/src/macros/common.tex.
% \providecommand keeps note-local definitions authoritative.
\usepackage{dsfont}
\providecommand{\ket}[1]{|#1\rangle}
\providecommand{\bra}[1]{\langle #1|}
\providecommand{\braket}[2]{\langle #1 | #2 \rangle}
\providecommand{\ketbra}[2]{|#1\rangle\!\langle #2|}
\providecommand{\Id}{\mathds{1}}
\providecommand{\one}{\mathds{1}}
\providecommand{\C}{\mathbb{C}}
\providecommand{\MN}[1]{M_{#1}(\C)}
\providecommand{\MD}{M_D(\C)}

% Verdict marker: \gapnote{<kind>}{<status>}. Kind is the policy
% classification (clarification, local-correction, scope-restriction,
% unfaithful, false-source, open-gap); status is open, wip (elimination
% underway), resolved, or historical. Machine-read by the site index;
% prints a verdict line.
\newcommand{\gapnote}[2]{\par\noindent\textbf{Verdict:} #1 (#2).\par}

\providecommand{\leanid}[1]{\textsf{\detokenize{#1}}}
\providecommand{\doclink}[1]{%
    \href{https://sirui-lu.com/TNLean/docs/#1.html}{\path{#1}}}
\providecommand{\gapnote}[2]{%
    \par\noindent\textbf{Verdict:} #1 (#2).\par}
\title{Automatic Numerical Admissibility of the Final Physical Scanner}
\author{}
\date{October 14, 2026}
\begin{document}
\maketitle
\gapnote{scope-restriction}{open}

The literal rounded final scales now imply the doubled-radius,
positive-separation and four-separation inequalities at one uniform
physical-size threshold.  The proof uses the accepted logarithmic-radius,
rounded-scale separation and positive band-count estimates.  It adds no
unproved numerical hypothesis.

The additive physical specialization
\leanid{TNLean.PEPS.AreaLaw.Scan.CollarScan.exists_sourceTarget_finalScale_mutualInformation_of_rounded_scales}
retains the same original-state information conclusions and every
geometric and entropy hypothesis of the conditional final-scale theorem.
Its radius coefficient precedes the source entropy coefficient.  The
information constant and enlarged threshold precede the domain, cut,
Hamiltonian, target and scanner.  The original conditional theorem remains
unchanged.

The remaining geometric inputs are the scanner with the actual domain
graph and ambient depth, support anchors, quadratic target cardinality,
ambient row counts, anchor multiplicity and cut clearance.  The genuine
source target and positive-prefix entropy estimates remain explicit.
Constructing these from the separated-template model and proving the
small safe-box exponent are separate obligations.  No complete
Proposition 9.6 or full area-law conclusion is asserted.

The numerical source is
\href{https://github.com/openai/math/blob/adc7f1241b42e322a6451854ab7e4b4c146bf78a/preprints/A-two-dimensional-area-law-from-a-global-spectral-gap-September-24-2026/build/sections/08-scanner.tex#L32-L58}{the scanner scale choices},
and the final application is
\href{https://github.com/openai/math/blob/adc7f1241b42e322a6451854ab7e4b4c146bf78a/preprints/A-two-dimensional-area-law-from-a-global-spectral-gap-September-24-2026/build/sections/08-scanner.tex#L808-L941}{the conditional collar-information argument}.
The broader work remains tracked by
\href{https://github.com/LionSR/TNLean/issues/8754}{\#8754},
\href{https://github.com/LionSR/TNLean/issues/8756}{\#8756}, and
\href{https://github.com/LionSR/TNLean/issues/9334}{\#9334}.

The proofs are independently written from the manuscript using existing
TNLean results.  No upstream Lean proof text is reused.
\end{document}
